\documentclass[11pt,a4paper]{article}

\usepackage[utf8]{inputenc}
\usepackage[T1]{fontenc}
\usepackage[french,provide=*]{babel}
\usepackage{amsmath,amssymb,amsthm,mathtools}
\usepackage{stmaryrd}
\usepackage{enumitem}
\usepackage[margin=2.4cm]{geometry}
\usepackage[expansion=false]{microtype}
\usepackage{hyperref}
\hypersetup{colorlinks=true,linkcolor=blue!55!black,citecolor=blue!55!black,urlcolor=blue!55!black}

% ------------------------------------------------------------------
% Environnements
% ------------------------------------------------------------------
\theoremstyle{plain}
\newtheorem{theorem}{Théorème}[section]
\newtheorem{proposition}[theorem]{Proposition}
\newtheorem{lemma}[theorem]{Lemme}
\newtheorem{corollary}[theorem]{Corollaire}

\theoremstyle{definition}
\newtheorem{definition}[theorem]{Définition}
\newtheorem{remark}[theorem]{Remarque}
\newtheorem{example}[theorem]{Exemple}

% ------------------------------------------------------------------
% Macros
% ------------------------------------------------------------------
\newcommand{\R}{\mathbb{R}}
\newcommand{\Q}{\mathbb{Q}}
\newcommand{\N}{\mathbb{N}}
\newcommand{\Rp}{\mathbb{R}_+}
\newcommand{\Rpb}{\overline{\mathbb{R}}_+}
\newcommand{\Rb}{\overline{\mathbb{R}}}
\newcommand{\B}{\mathcal{B}}
\newcommand{\F}{\mathcal{F}}
\newcommand{\A}{\mathcal{A}}
\newcommand{\Bo}{\mathcal{B}}
\newcommand{\sig}{\sigma}
\newcommand{\ind}{\mathbf{1}}
\newcommand{\eps}{\varepsilon}
\newcommand{\ra}{\rightarrow}
\newcommand{\limsupn}{\limsup_{n\to\infty}}
\newcommand{\liminfn}{\liminf_{n\to\infty}}
\DeclareMathOperator{\Vois}{\mathcal{V}}
\newcommand{\restr}[2]{{#1}\big|_{#2}}

\title{\textbf{Extension de la mesurabilité à $\Rpb=[0,+\infty]$}\\[4pt]
\large Compactifications, topologie de l'ordre et tribu borélienne}
\author{}
\date{\today}

\begin{document}
\maketitle

\begin{abstract}
On construit rigoureusement la droite positive achevée $\Rpb=[0,+\infty]$ comme
espace topologique compact métrisable, en montrant la coïncidence de trois
compactifications naturelles (topologie de l'ordre, compactifié d'Alexandroff,
plongement métrique dans $[0,1]$). On en déduit la structure de la tribu
borélienne $\B(\Rpb)$, ses familles génératrices, la caractérisation par trace,
puis on développe la théorie de la mesurabilité des fonctions à valeurs dans
$\Rpb$. On démontre en détail la stabilité par $\sup$, $\inf$, $\limsup$,
$\liminf$, la mesurabilité des opérations arithmétiques étendues (addition,
multiplication avec la convention $0\cdot\infty=0$), ainsi que le principe de
transfert par homéomorphisme. Toutes les propriétés topologiques et
ensemblistes utilisées sont énoncées et prouvées.
\end{abstract}

\tableofcontents

% ==================================================================
\section{Motivation : pourquoi ajouter le point $+\infty$}
% ==================================================================

En théorie de la mesure et des probabilités, les grandeurs fondamentales
(mesures d'ensembles, intégrales de fonctions positives, sommes de séries à
termes positifs) prennent naturellement leurs valeurs dans $[0,+\infty]$ et non
dans $[0,+\infty)$. Deux raisons structurelles imposent l'adjonction du point
$+\infty$.

\paragraph{(i) Complétude du treillis ordonné.}
On rappelle qu'un ensemble totalement ordonné $(L,\le)$ est un \emph{treillis
complet} si toute partie $A\subseteq L$ admet une borne supérieure $\sup A$ et
une borne inférieure $\inf A$ dans $L$. L'ensemble $\Rp=[0,+\infty)$ n'est
\emph{pas} un treillis complet : la partie $A=\N$ n'y possède pas de borne
supérieure. En revanche $\Rpb=[0,+\infty]$, muni de l'ordre prolongé par
$x\le+\infty$ pour tout $x$, est un treillis complet (Proposition~\ref{prop:treillis}).
C'est cette complétude qui garantit l'existence de $\sup_n f_n$ et
$\inf_n f_n$ pour une suite quelconque de fonctions, donc celle de
$\limsup$ et $\liminf$.

\paragraph{(ii) Compacité.}
Le point $+\infty$ n'est pas un simple symbole formel : c'est un point d'un
espace topologique \emph{compact métrisable} homéomorphe au segment $[0,1]$.
Cette compacité fournit l'existence de valeurs d'adhérence pour toute suite,
la régularité des mesures boréliennes finies, et le cadre d'espace borélien
standard dans lequel les variables aléatoires se comportent le mieux.

L'objet de ce texte est de rendre ces deux points parfaitement explicites et
d'en déduire toute la théorie de la mesurabilité étendue.

% ==================================================================
\section{Rappels d'ordre et structure de treillis complet}
% ==================================================================

\begin{definition}[Droite positive achevée]
On pose $\Rpb=\Rp\cup\{+\infty\}=[0,+\infty]$, où $+\infty$ est un élément
n'appartenant pas à $\Rp$. On prolonge l'ordre total $\le$ de $\Rp$ en posant
$x\le+\infty$ pour tout $x\in\Rpb$. Alors $(\Rpb,\le)$ est un ensemble
totalement ordonné admettant un plus petit élément $0$ et un plus grand
élément $+\infty$.
\end{definition}

\begin{proposition}[$\Rpb$ est un treillis complet]\label{prop:treillis}
Toute partie $A\subseteq\Rpb$ admet une borne supérieure et une borne
inférieure dans $\Rpb$. Avec les conventions $\sup\varnothing=0$ et
$\inf\varnothing=+\infty$.
\end{proposition}

\begin{proof}
Traitons la borne supérieure ; le cas de la borne inférieure est symétrique
(ou s'obtient via $\inf A=\sup\{m\in\Rpb: \forall a\in A,\ m\le a\}$).

Soit $A\subseteq\Rpb$. Distinguons trois cas.

\emph{Cas 1 : $A=\varnothing$.} Tout élément de $\Rpb$ est un majorant (vide de
condition) ; le plus petit majorant est le plus petit élément $0$. Donc
$\sup\varnothing=0$.

\emph{Cas 2 : $+\infty\in A$, ou $A\cap\Rp$ n'est pas majorée dans $\Rp$.}
Alors $+\infty$ est un majorant de $A$, et c'est le plus grand élément de
$\Rpb$, donc le plus petit majorant possible dès qu'aucun réel ne majore $A$.
Si $+\infty\in A$, aucun élément $<+\infty$ ne majore $A$ ; si $A\cap\Rp$ est
non bornée, aucun réel ne la majore. Dans les deux cas $\sup A=+\infty$.

\emph{Cas 3 : $A\subseteq\Rp$ et $A$ est majorée dans $\Rp$ (i.e. $\exists
M\in\Rp,\ \forall a\in A,\ a\le M$).} Si $A=\varnothing$ on est dans le Cas~1 ;
sinon $A$ est une partie non vide majorée de $\R$, donc par la propriété de la
borne supérieure de $\R$ (axiome de complétude de l'ordre de $\R$), $\sup A$
existe dans $\R$ et est $\ge0$, donc appartient à $\Rp\subseteq\Rpb$.

Dans tous les cas $\sup A$ existe dans $\Rpb$.
\end{proof}

\begin{remark}
La Proposition~\ref{prop:treillis} est le socle \emph{algébrique} de toute la
suite : elle assure que les opérations $\sup$ et $\inf$ ne « sortent » jamais de
$\Rpb$. La compacité topologique établie plus bas en est la contrepartie
\emph{analytique}.
\end{remark}

% ==================================================================
\section{Les trois compactifications de $\Rp$}
% ==================================================================

On munit $\Rpb$ d'une topologie et l'on montre que trois constructions
naturelles donnent la même topologie. On note $\Vois(x)$ le filtre des
voisinages d'un point $x$.

% ------------------------------------------------------------------
\subsection{Topologie de l'ordre}
% ------------------------------------------------------------------

\begin{definition}[Topologie de l'ordre]\label{def:ordre}
Sur un ensemble totalement ordonné $(L,\le)$, la \emph{topologie de l'ordre}
$\tau_{\le}$ est la topologie engendrée par la prébase des rayons ouverts
\[
\{x\in L: x<a\}\quad(a\in L),\qquad \{x\in L: x>a\}\quad(a\in L).
\]
\end{definition}

Pour $L=\Rpb$, décrivons une base explicite. Les rayons sont
$[0,a)=\{x<a\}$ pour $a\in\Rpb$ et $(a,+\infty]=\{x>a\}$ pour $a\in\Rpb$
(en convenant $[0,0)=\varnothing$, $(+\infty,+\infty]=\varnothing$). Les
intersections finies de tels rayons donnent la base
\begin{equation}\label{eq:base-ordre}
\mathcal{B}_{\le}=\big\{[0,b): b\in\Rpb\big\}\;\cup\;
\big\{(a,b): 0\le a<b\le+\infty\big\}\;\cup\;
\big\{(a,+\infty]: a\in\Rp\big\}.
\end{equation}

\begin{proposition}[Voisinages de $+\infty$ et de $0$]\label{prop:vois}
Dans $(\Rpb,\tau_{\le})$ :
\begin{enumerate}[label=(\alph*)]
\item les ensembles $(a,+\infty]$, $a\in\Rp$, forment une base de voisinages
de $+\infty$ ;
\item les ensembles $[0,b)$, $b\in\Rp^*$, forment une base de voisinages
de $0$ ;
\item la topologie induite par $\tau_{\le}$ sur $\Rp$ coïncide avec la
topologie usuelle de $\Rp$.
\end{enumerate}
\end{proposition}

\begin{proof}
(a) Un élément de base contenant $+\infty$ doit être de la forme $(a,+\infty]$
d'après~\eqref{eq:base-ordre} (les $[0,b)$ et $(a,b)$ ne contiennent pas
$+\infty$). Réciproquement chaque $(a,+\infty]$ est ouvert et contient
$+\infty$. D'où le résultat. (b) Symétrique. (c) Pour $a,b\in\Rp$ les traces
sur $\Rp$ des éléments de base sont $[0,b)$, $(a,b)$ et $(a,+\infty)$, qui
engendrent exactement la topologie usuelle de $[0,+\infty)$.
\end{proof}

% ------------------------------------------------------------------
\subsection{Plongement métrique dans $[0,1]$}
% ------------------------------------------------------------------

\begin{definition}[Homéomorphisme de compactification]\label{def:phi}
On définit $\varphi:\Rpb\to[0,1]$ par
\[
\varphi(x)=\frac{x}{1+x}\quad(x\in\Rp),\qquad \varphi(+\infty)=1.
\]
\end{definition}

\begin{proposition}[$\varphi$ est une bijection croissante]\label{prop:phi-bij}
L'application $\varphi$ est une bijection strictement croissante de $\Rpb$ sur
$[0,1]$, de réciproque
\[
\varphi^{-1}(y)=\frac{y}{1-y}\quad(y\in[0,1)),\qquad \varphi^{-1}(1)=+\infty.
\]
\end{proposition}

\begin{proof}
Sur $\Rp$, $\varphi(x)=1-\frac{1}{1+x}$ est dérivable de dérivée
$\varphi'(x)=(1+x)^{-2}>0$, donc strictement croissante ; elle réalise une
bijection de $[0,+\infty)$ sur $[0,1)$ (limites $\varphi(0)=0$,
$\varphi(x)\to1^-$). En posant $\varphi(+\infty)=1$ on prolonge en une
bijection croissante $\Rpb\to[0,1]$ car $x<+\infty\Rightarrow\varphi(x)<1$.
L'expression de $\varphi^{-1}$ se vérifie par calcul direct : si
$y=\frac{x}{1+x}$ alors $x=\frac{y}{1-y}$.
\end{proof}

\begin{definition}[Métrique de $\Rpb$]\label{def:metrique}
On pose, pour $x,y\in\Rpb$,
\[
d(x,y)=\big|\varphi(x)-\varphi(y)\big|.
\]
C'est une distance sur $\Rpb$ (image réciproque de la distance usuelle de
$[0,1]$ par l'injection $\varphi$).
\end{definition}

\begin{proposition}[La topologie métrique coïncide avec $\tau_\le$]\label{prop:coinc-metrique}
La topologie $\tau_d$ induite par $d$ sur $\Rpb$ est égale à la topologie de
l'ordre $\tau_\le$. De plus $\varphi:(\Rpb,\tau_\le)\to([0,1],|\cdot|)$ est un
homéomorphisme.
\end{proposition}

\begin{proof}
Par construction $\varphi$ est une isométrie de $(\Rpb,d)$ sur
$([0,1],|\cdot|)$, donc un homéomorphisme pour $\tau_d$. Il reste à voir que
$\varphi$ est aussi un homéomorphisme pour $\tau_\le$, ce qui entraînera
$\tau_d=\tau_\le$ (deux topologies rendant la même bijection homéomorphe vers
le même espace coïncident).

Comme $\varphi$ est une bijection croissante entre ensembles totalement
ordonnés, elle transporte les rayons sur les rayons :
$\varphi(\{x<a\})=\{y<\varphi(a)\}\cap[0,1]$ et de même pour $>$. Elle
transforme donc la prébase de la topologie de l'ordre de $\Rpb$ en la prébase
de la topologie de l'ordre de $[0,1]$. Or sur $[0,1]$ la topologie de l'ordre
est la topologie usuelle. Donc $\varphi$ est un homéomorphisme
$(\Rpb,\tau_\le)\to[0,1]$.
\end{proof}

\begin{corollary}[Compacité et métrisabilité]\label{cor:compact}
$(\Rpb,\tau_\le)$ est un espace \emph{compact}, \emph{métrisable}, à base
dénombrable, séparable et complet. En particulier c'est un espace
\emph{polonais compact}.
\end{corollary}

\begin{proof}
$[0,1]$ est compact (Heine--Borel), métrique complet, séparable. Ces
propriétés sont invariantes par homéomorphisme (la métrisabilité, la
séparabilité, la base dénombrable et la compacité le sont ; la complétude est
transportée par l'isométrie $\varphi$). Le résultat suit de
Proposition~\ref{prop:coinc-metrique}.
\end{proof}

% ------------------------------------------------------------------
\subsection{Compactifié d'Alexandroff}
% ------------------------------------------------------------------

\begin{definition}[Compactifié d'Alexandroff]\label{def:alex}
Soit $X$ un espace topologique localement compact séparé non compact. Le
\emph{compactifié d'Alexandroff} (ou compactifié en un point) est
$X^*=X\cup\{\infty\}$ muni de la topologie
\[
\tau^*=\tau_X\;\cup\;\big\{\{\infty\}\cup(X\setminus K): K\subseteq X\text{ compact}\big\},
\]
où $\tau_X$ est la topologie de $X$. Les ouverts de $X^*$ sont donc les ouverts
de $X$ et les complémentaires (dans $X^*$) des compacts de $X$.
\end{definition}

On rappelle sans démonstration le résultat classique : si $X$ est localement
compact séparé non compact, alors $X^*$ est compact séparé et $X$ en est un
sous-espace ouvert dense.

\begin{proposition}[Alexandroff $=$ ordre pour $\Rp$]\label{prop:alex-ordre}
$\Rp=[0,+\infty)$ est localement compact, séparé, non compact. Son compactifié
d'Alexandroff $\Rp^*=\Rp\cup\{\infty\}$ est homéomorphe à $(\Rpb,\tau_\le)$ via
l'identification $\infty\leftrightarrow+\infty$.
\end{proposition}

\begin{proof}
$\Rp$ est un sous-espace de $\R$, séparé ; il est localement compact car tout
point $x$ possède le voisinage compact $[\max(0,x-1),x+1]$ ; il n'est pas
compact car non borné (ou : le recouvrement $\{[0,n)\}_n$ n'a pas de
sous-recouvrement fini). Le compactifié d'Alexandroff est donc bien défini.

Identifions $\infty$ avec $+\infty$, de sorte que $\Rp^*=\Rpb$ en tant
qu'ensembles. Il reste à comparer les topologies. Sur $\Rp$, $\tau^*$ et
$\tau_\le$ induisent toutes deux la topologie usuelle (Def.~\ref{def:alex}
pour $\tau^*$ ; Proposition~\ref{prop:vois}(c) pour $\tau_\le$). Il suffit donc
de comparer les voisinages du point ajouté.

Par le théorème de Heine--Borel dans $\Rp$, les parties compactes de $\Rp$ sont
exactement les fermés bornés. Un voisinage de base de $\infty$ dans $\tau^*$ est
$\{\infty\}\cup(\Rp\setminus K)$ avec $K$ compact. Comme $K$ est borné, il
existe $M$ tel que $K\subseteq[0,M]$, d'où
$\{\infty\}\cup(\Rp\setminus K)\supseteq\{\infty\}\cup(M,+\infty)=(M,+\infty]$.
Réciproquement $(M,+\infty]=\{\infty\}\cup(\Rp\setminus[0,M])$ est un tel
voisinage (avec $K=[0,M]$ compact). Les voisinages de $\infty$ dans $\tau^*$ ont
donc pour base les $(M,+\infty]$, $M\in\Rp$, qui sont exactement les voisinages
de base de $+\infty$ dans $\tau_\le$ (Proposition~\ref{prop:vois}(a)). Les deux
topologies coïncident.
\end{proof}

% ------------------------------------------------------------------
\subsection{Interprétation par les bouts et comparaison avec $\R$}
% ------------------------------------------------------------------

\begin{remark}[Nombre de bouts]
Intuitivement, le nombre de points à ajouter pour compactifier « correctement »
(en préservant l'ordre) est le nombre de \emph{bouts} de l'espace au sens de
Freudenthal.
\begin{itemize}[nosep]
\item $\Rp=[0,+\infty)$ possède un \emph{unique bout} (à l'infini ; le bord $0$
est déjà présent). Une seule extension raisonnable existe : le point $+\infty$.
Les trois compactifications coïncident, et l'on tombe sur le segment $[0,1]$.
\item $\R$ possède \emph{deux bouts} ($-\infty$ et $+\infty$). La droite achevée
$\Rb=[-\infty,+\infty]$ des probabilistes est la compactification à
\emph{deux points}, homéomorphe à $[-1,1]$ ; ce n'est pas le compactifié
d'Alexandroff de $\R$.
\item Le compactifié d'Alexandroff de $\R$ (un seul point recollant $\pm\infty$)
est homéomorphe au cercle $S^1$ ; il détruit l'ordre et est donc inadapté à
l'intégration.
\end{itemize}
Ainsi la situation de $\Rp$ est plus simple que celle de $\R$ : coïncidence
totale des compactifications.
\end{remark}

% ==================================================================
\section{Base dénombrable et structure topologique fine}
% ==================================================================

\begin{proposition}[Base dénombrable explicite]\label{prop:base-denombrable}
La famille dénombrable
\[
\mathcal{B}_0=\big\{[0,q): q\in\Q_+^*\big\}\;\cup\;
\big\{(p,q): p,q\in\Q_+,\ p<q\big\}\;\cup\;
\big\{(p,+\infty]: p\in\Q_+\big\}
\]
est une base de la topologie de $\Rpb$.
\end{proposition}

\begin{proof}
Chaque élément de $\mathcal{B}_0$ est ouvert. Montrons que tout ouvert est
réunion d'éléments de $\mathcal{B}_0$, i.e. que pour tout ouvert $U$ et tout
$x\in U$ il existe $B\in\mathcal{B}_0$ avec $x\in B\subseteq U$.

Comme $\mathcal{B}_\le$ de~\eqref{eq:base-ordre} est une base, il suffit de
raffiner chaque type d'élément par des bornes rationnelles.
\begin{itemize}[nosep]
\item Si $x\in(a,b)$ avec $a,b\in\Rp$ : par densité de $\Q$ dans $\R$ on choisit
$p,q\in\Q_+$ avec $a<p<x<q<b$, d'où $x\in(p,q)\subseteq(a,b)$.
\item Si $x\in[0,b)$ : on choisit $q\in\Q_+^*$ avec $x<q<b$ (si $x=0$) ou
$p,q\in\Q$ avec $\ldots$ ; plus simplement $x\in[0,q)\subseteq[0,b)$ pour
$q\in(x,b)\cap\Q$ si $x=0$, et sinon on se ramène au cas $(p,q)$.
\item Si $x\in(a,+\infty]$ : on choisit $p\in\Q_+$ avec $a<p<x$ (possible car
$x>a\ge0$, en prenant $p<x$ fini si $x<\infty$, et n'importe quel $p>a$ si
$x=+\infty$), d'où $x\in(p,+\infty]\subseteq(a,+\infty]$.
\end{itemize}
Dans tous les cas on a intercalé un élément de $\mathcal{B}_0$. Donc
$\mathcal{B}_0$ est une base.
\end{proof}

\begin{corollary}[Espace borélien standard]\label{cor:standard}
$(\Rpb,\B(\Rpb))$ est un \emph{espace borélien standard} : il existe un
isomorphisme borélien entre $(\Rpb,\B(\Rpb))$ et $([0,1],\B([0,1]))$, et par le
théorème d'isomorphisme de Kuratowski il est borélo-isomorphe à
$(\R,\B(\R))$.
\end{corollary}

\begin{proof}
$\Rpb$ est polonais compact (Corollaire~\ref{cor:compact}), donc son espace
borélien est standard. L'homéomorphisme $\varphi$ (Def.~\ref{def:phi}) est en
particulier un isomorphisme borélien $\Rpb\to[0,1]$. Le théorème d'isomorphisme
de Kuratowski affirme que deux espaces boréliens standards de même cardinalité
du continu sont borélo-isomorphes.
\end{proof}

% ==================================================================
\section{La tribu borélienne $\B(\Rpb)$}
% ==================================================================

\begin{definition}
$\B(\Rpb)=\sig(\tau_\le)$ est la tribu engendrée par les ouverts de $\Rpb$.
\end{definition}

% ------------------------------------------------------------------
\subsection{Familles génératrices}
% ------------------------------------------------------------------

\begin{theorem}[Générateurs de $\B(\Rpb)$]\label{thm:generateurs}
On a les égalités de tribus
\[
\B(\Rpb)
=\sig\big(\{(a,+\infty]: a\in\Rp\}\big)
=\sig\big(\{[0,a]: a\in\Rp\}\big)
=\sig\big(\{[a,+\infty]: a\in\Rp\}\big)
=\sig\big(\{[0,a): a\in\Rp\}\big).
\]
De plus on peut restreindre $a$ à $\Q_+$ dans chacune de ces familles.
\end{theorem}

\begin{proof}
Notons $\mathcal{G}_1=\{(a,+\infty]:a\in\Rp\}$ et $\sig(\mathcal{G}_1)$ la tribu
qu'elle engendre.

\emph{$\sig(\mathcal{G}_1)\subseteq\B(\Rpb)$ :} chaque $(a,+\infty]$ est ouvert
(Proposition~\ref{prop:vois}), donc borélien.

\emph{$\B(\Rpb)\subseteq\sig(\mathcal{G}_1)$ :} il suffit de montrer que tout
élément de la base dénombrable $\mathcal{B}_0$
(Proposition~\ref{prop:base-denombrable}) appartient à $\sig(\mathcal{G}_1)$,
car alors tout ouvert, réunion \emph{dénombrable} d'éléments de $\mathcal{B}_0$,
y appartient aussi. Or, pour $p,q\in\Q_+$ :
\begin{align*}
[0,q]&=\Rpb\setminus(q,+\infty]\in\sig(\mathcal{G}_1),\\
[0,q)&=\bigcup_{n\ge1}\Big[0,q-\tfrac1n\Big]\in\sig(\mathcal{G}_1)
\quad(\text{réunion dénombrable}),\\
(p,q)&=(p,+\infty]\cap[0,q)\in\sig(\mathcal{G}_1),\\
(p,+\infty]&\in\mathcal{G}_1\subseteq\sig(\mathcal{G}_1).
\end{align*}
Donc $\mathcal{B}_0\subseteq\sig(\mathcal{G}_1)$, puis
$\B(\Rpb)\subseteq\sig(\mathcal{G}_1)$. D'où l'égalité.

\emph{Autres familles.} Chacune s'obtient de $\mathcal{G}_1$ par
complémentation ou réunion/intersection dénombrable :
\[
[0,a]=\Rpb\setminus(a,+\infty],\qquad
[a,+\infty]=\bigcap_{n\ge1}\Big(a-\tfrac1n,+\infty\Big],\qquad
[0,a)=\bigcup_{n\ge1}\Big[0,a-\tfrac1n\Big],
\]
et réciproquement $(a,+\infty]=\Rpb\setminus[0,a]$, ce qui montre que les
quatre tribus engendrées coïncident. Enfin la restriction à $a\in\Q_+$ est
licite car tout $(a,+\infty]$ avec $a\in\Rp$ s'écrit
$\bigcup_{q\in\Q,\,q>a}(q,+\infty]$ (réunion dénombrable), et de même pour les
autres familles par densité de $\Q$.
\end{proof}

% ------------------------------------------------------------------
\subsection{Le singleton $\{+\infty\}$ et la caractérisation par trace}
% ------------------------------------------------------------------

\begin{lemma}[$\{+\infty\}$ est un fermé $G_\delta$ borélien]\label{lem:infty-borel}
Le singleton $\{+\infty\}$ est fermé dans $\Rpb$, s'écrit comme intersection
dénombrable d'ouverts, et est borélien. De même $\Rp=[0,+\infty)$ est ouvert,
$F_\sigma$, et borélien.
\end{lemma}

\begin{proof}
On a
\[
\{+\infty\}=\bigcap_{n\ge1}(n,+\infty],
\]
intersection dénombrable d'ouverts : donc $\{+\infty\}$ est un $G_\delta$, en
particulier borélien. Son complémentaire
$\Rp=[0,+\infty)=\bigcup_{n\ge1}[0,n]=\bigcup_{n\ge1}[0,n)$ est une réunion
dénombrable de fermés (resp.\ d'ouverts), donc un ouvert $F_\sigma$ borélien.
(La fermeture de $\{+\infty\}$ résulte aussi de la séparation de $\Rpb$.)
\end{proof}

\begin{theorem}[Caractérisation par trace]\label{thm:trace}
On a l'égalité
\[
\B(\Rpb)=\Big\{A\subseteq\Rpb:\ A\cap\Rp\in\B(\Rp)\Big\},
\]
et de façon équivalente
\[
\B(\Rpb)=\Big\{B:\ B\in\B(\Rp)\Big\}\;\cup\;\Big\{B\cup\{+\infty\}:\ B\in\B(\Rp)\Big\}.
\]
Autrement dit, un borélien de $[0,+\infty]$ est exactement un borélien de
$[0,+\infty)$, auquel on adjoint éventuellement le point $+\infty$.
Réciproquement, $\B(\Rp)$ est la tribu trace : $\B(\Rp)=\B(\Rpb)\cap\Rp
:=\{A\cap\Rp:A\in\B(\Rpb)\}$.
\end{theorem}

\begin{proof}
\emph{Trace.} Rappelons le fait général : si $(X,\tau)$ est un espace
topologique et $Y\subseteq X$, alors $\B(Y)=\B(X)\cap Y$, où la topologie de $Y$
est la topologie trace (la tribu borélienne de la trace est la trace de la
tribu borélienne). Appliqué à $X=\Rpb$, $Y=\Rp$ (dont la topologie trace est
l'usuelle par Proposition~\ref{prop:vois}(c)), cela donne
$\B(\Rp)=\B(\Rpb)\cap\Rp$. Prouvons ce fait.

Posons $\mathcal{T}=\{A\cap Y:A\in\B(X)\}$. C'est une tribu sur $Y$ (image
réciproque d'une tribu par l'inclusion $Y\hookrightarrow X$, qui est
$\{A\cap Y\}$). Elle contient les ouverts de $Y$ (traces des ouverts de $X$),
donc $\B(Y)\subseteq\mathcal{T}$. Réciproquement, la famille
$\mathcal{S}=\{A\subseteq X: A\cap Y\in\B(Y)\}$ est une tribu sur $X$ (car
$A\mapsto A\cap Y$ commute au complémentaire et aux réunions dénombrables)
contenant les ouverts de $X$, donc $\B(X)\subseteq\mathcal{S}$, ce qui signifie
$A\cap Y\in\B(Y)$ pour tout $A\in\B(X)$, i.e.\ $\mathcal{T}\subseteq\B(Y)$. D'où
$\mathcal{T}=\B(Y)$.

\emph{Description explicite.} Soit $A\subseteq\Rpb$ tel que
$B:=A\cap\Rp\in\B(\Rp)$. D'après ce qui précède, $\B(\Rp)=\B(\Rpb)\cap\Rp$, donc
il existe $A'\in\B(\Rpb)$ avec $B=A'\cap\Rp$. Deux cas :
\begin{itemize}[nosep]
\item si $+\infty\notin A$, alors $A=B$. Or $B\in\B(\Rp)\subseteq\B(\Rpb)$ car
$\B(\Rp)=\B(\Rpb)\cap\Rp\subseteq\B(\Rpb)$ (tout élément de la trace est de la
forme $A'\cap\Rp$, et $\Rp\in\B(\Rpb)$ par Lemme~\ref{lem:infty-borel}, donc
$A'\cap\Rp\in\B(\Rpb)$). Ainsi $A\in\B(\Rpb)$ ;
\item si $+\infty\in A$, alors $A=B\cup\{+\infty\}$ avec $B\in\B(\Rpb)$ et
$\{+\infty\}\in\B(\Rpb)$ (Lemme~\ref{lem:infty-borel}), donc
$A\in\B(\Rpb)$.
\end{itemize}
Réciproquement, si $A\in\B(\Rpb)$, alors $A\cap\Rp\in\B(\Rpb)\cap\Rp=\B(\Rp)$.
Ceci établit la double égalité annoncée.
\end{proof}

\begin{remark}
La caractérisation par trace est la description la plus commode en pratique :
pour vérifier qu'une partie de $[0,+\infty]$ est borélienne, il suffit de
regarder sa trace sur $[0,+\infty)$ (calcul « fini ») et de traiter à part le
point $+\infty$, qui est toujours borélien.
\end{remark}

% ==================================================================
\section{Fonctions mesurables à valeurs dans $\Rpb$}
% ==================================================================

Dans toute cette section, $(\Omega,\F)$ désigne un espace mesurable.

\begin{definition}[Mesurabilité]
Une application $f:\Omega\to\Rpb$ est \emph{$(\F,\B(\Rpb))$-mesurable} (en abrégé
\emph{mesurable}) si $f^{-1}(A)\in\F$ pour tout $A\in\B(\Rpb)$.
\end{definition}

% ------------------------------------------------------------------
\subsection{Critères de mesurabilité}
% ------------------------------------------------------------------

\begin{proposition}[Réduction à une famille génératrice]\label{prop:crit-gen}
Soit $\mathcal{G}$ une famille engendrant $\B(\Rpb)$, i.e.\
$\sig(\mathcal{G})=\B(\Rpb)$. Alors $f:\Omega\to\Rpb$ est mesurable si et
seulement si $f^{-1}(G)\in\F$ pour tout $G\in\mathcal{G}$.
\end{proposition}

\begin{proof}
Le sens direct est clair. Réciproquement, supposons $f^{-1}(G)\in\F$ pour tout
$G\in\mathcal{G}$. La famille $\mathcal{M}=\{A\subseteq\Rpb: f^{-1}(A)\in\F\}$
est une tribu, car $f^{-1}$ commute au passage au complémentaire et aux
réunions dénombrables :
\[
f^{-1}(A^c)=f^{-1}(A)^c,\qquad
f^{-1}\Big(\bigcup_n A_n\Big)=\bigcup_n f^{-1}(A_n).
\]
Par hypothèse $\mathcal{G}\subseteq\mathcal{M}$, donc
$\B(\Rpb)=\sig(\mathcal{G})\subseteq\mathcal{M}$, ce qui signifie que $f$ est
mesurable.
\end{proof}

\begin{corollary}[Critères usuels]\label{cor:criteres}
Pour $f:\Omega\to\Rpb$, les assertions suivantes sont équivalentes :
\begin{enumerate}[label=(\roman*),nosep]
\item $f$ est mesurable ;
\item $\{f>a\}\in\F$ pour tout $a\in\Rp$ (resp.\ pour tout $a\in\Q_+$) ;
\item $\{f\ge a\}\in\F$ pour tout $a\in\Rp$ ;
\item $\{f\le a\}\in\F$ pour tout $a\in\Rp$ ;
\item $\{f<a\}\in\F$ pour tout $a\in\Rp$.
\end{enumerate}
Ici $\{f>a\}:=f^{-1}((a,+\infty])$, etc.
\end{corollary}

\begin{proof}
Application directe de la Proposition~\ref{prop:crit-gen} aux familles
génératrices du Théorème~\ref{thm:generateurs} :
$\{(a,+\infty]\}$, $\{[a,+\infty]\}$, $\{[0,a]\}$, $\{[0,a)\}$.
\end{proof}

\begin{proposition}[Gestion automatique de l'infini]\label{prop:infini-auto}
Si $f:\Omega\to\Rpb$ est mesurable, alors
\[
\{f=+\infty\}=\bigcap_{n\ge1}\{f>n\}\in\F,\qquad
\{f<+\infty\}=\bigcup_{n\ge1}\{f\le n\}\in\F.
\]
\end{proposition}

\begin{proof}
On a $f(\omega)=+\infty$ ssi $f(\omega)>n$ pour tout $n$, d'où la première
égalité ensembliste ; les $\{f>n\}$ sont dans $\F$ par
Corollaire~\ref{cor:criteres}, et $\F$ est stable par intersection
dénombrable. La seconde égalité est le complémentaire.
\end{proof}

% ------------------------------------------------------------------
\subsection{Stabilité par les opérations de treillis}
% ------------------------------------------------------------------

C'est le c\oe ur de l'intérêt de $\Rpb$ : la complétude du treillis
(Proposition~\ref{prop:treillis}) se traduit en stabilité de la mesurabilité.

\begin{theorem}[Sup et inf dénombrables]\label{thm:supinf}
Soit $(f_n)_{n\ge1}$ une suite de fonctions mesurables $\Omega\to\Rpb$. Alors
les fonctions
\[
\Big(\sup_n f_n\Big)(\omega)=\sup_n f_n(\omega),\qquad
\Big(\inf_n f_n\Big)(\omega)=\inf_n f_n(\omega)
\]
sont bien définies (à valeurs dans $\Rpb$, par
Proposition~\ref{prop:treillis}) et mesurables.
\end{theorem}

\begin{proof}
Posons $g=\sup_n f_n$. Pour tout $a\in\Rp$, l'équivalence
\[
g(\omega)\le a\iff \forall n,\ f_n(\omega)\le a
\]
(caractérisation de la borne supérieure comme plus petit majorant) donne
l'égalité ensembliste
\[
\{g\le a\}=\bigcap_{n\ge1}\{f_n\le a\}.
\]
Chaque $\{f_n\le a\}\in\F$ (Corollaire~\ref{cor:criteres}) et $\F$ est stable
par intersection dénombrable, donc $\{g\le a\}\in\F$ pour tout $a\in\Rp$. Par
Corollaire~\ref{cor:criteres}(iv), $g$ est mesurable.

Symétriquement, pour $h=\inf_n f_n$ et tout $a\in\Rp$,
\[
h(\omega)\ge a\iff\forall n,\ f_n(\omega)\ge a
\quad\Longrightarrow\quad
\{h\ge a\}=\bigcap_{n\ge1}\{f_n\ge a\}\in\F,
\]
d'où la mesurabilité de $h$ par Corollaire~\ref{cor:criteres}(iii).
\end{proof}

\begin{theorem}[Limsup, liminf, limite]\label{thm:limsup}
Soit $(f_n)$ une suite de fonctions mesurables $\Omega\to\Rpb$. Les fonctions
\[
\limsupn f_n=\inf_{n\ge1}\ \sup_{k\ge n} f_k,\qquad
\liminfn f_n=\sup_{n\ge1}\ \inf_{k\ge n} f_k
\]
sont bien définies et mesurables, à valeurs dans $\Rpb$. En particulier
l'ensemble de convergence
\[
C=\Big\{\omega:\ \lim_n f_n(\omega)\text{ existe dans }\Rpb\Big\}
=\Big\{\limsupn f_n=\liminfn f_n\Big\}
\]
est mesurable, et sur $C$ la limite $\lim_n f_n$ est mesurable.
\end{theorem}

\begin{proof}
Pour chaque $n$, $g_n:=\sup_{k\ge n} f_k$ est mesurable
(Théorème~\ref{thm:supinf} appliqué à la sous-suite $(f_k)_{k\ge n}$). La suite
$(g_n)$ est décroissante ; $\limsup_n f_n=\inf_n g_n$ est mesurable
(Théorème~\ref{thm:supinf} à nouveau). De même
$h_n:=\inf_{k\ge n}f_k$ est mesurable et $\liminf_n f_n=\sup_n h_n$ l'est.

L'existence dans $\Rpb$ est garantie par la complétude du treillis
(Proposition~\ref{prop:treillis}) : les bornes emboîtées existent toujours.
Enfin, comme $\liminf\le\limsup$ toujours, l'ensemble de convergence est
\[
C=\{\limsup_n f_n=\liminf_n f_n\}
=\{\limsup_n f_n\le\liminf_n f_n\},
\]
qui est mesurable car différence/égalité de deux fonctions mesurables à valeurs
dans $\Rpb$ : plus précisément, pour deux fonctions mesurables $u,v:\Omega\to\Rpb$,
\[
\{u\le v\}=\bigcap_{q\in\Q_+}\big(\{u\le q\}\cup\{v> q\}\big)\ \cup\ \{v=+\infty\}
\in\F,
\]
(voir Lemme~\ref{lem:comparaison} ci-dessous). Sur $C$, $\lim_n f_n$ coïncide
avec $\limsup_n f_n$, donc est mesurable.
\end{proof}

\begin{lemma}[Mesurabilité des comparaisons]\label{lem:comparaison}
Si $u,v:\Omega\to\Rpb$ sont mesurables, alors $\{u\le v\}$, $\{u<v\}$,
$\{u=v\}$ appartiennent à $\F$.
\end{lemma}

\begin{proof}
On utilise la densité de $\Q_+$ et la relation, valable dans le treillis
totalement ordonné $\Rpb$ :
\[
u(\omega)<v(\omega)\iff \exists q\in\Q_+,\ u(\omega)<q<v(\omega)\ \text{ ou }\ 
\big(u(\omega)<+\infty=v(\omega)\ \text{avec intercalation}\big).
\]
Plus proprement, pour $x,y\in\Rpb$ on a $x<y$ si et seulement s'il existe
$q\in\Q_+$ tel que $x<q\le y$ lorsque $y<+\infty$, et $x<y=+\infty$ équivaut à
$x<+\infty$. On écrit donc
\[
\{u<v\}=\Big(\bigcup_{q\in\Q_+}\{u<q\}\cap\{v>q\}\Big)\ \cup\ 
\big(\{u<+\infty\}\cap\{v=+\infty\}\big).
\]
Chaque terme est dans $\F$ (Corollaire~\ref{cor:criteres} et
Proposition~\ref{prop:infini-auto}), donc $\{u<v\}\in\F$. Alors
$\{u\ge v\}=\{u<v\}^c\in\F$, $\{v<u\}\in\F$ par symétrie,
$\{u=v\}=\{u<v\}^c\cap\{v<u\}^c\in\F$ et $\{u\le v\}=\{v<u\}^c\in\F$.
\end{proof}

% ------------------------------------------------------------------
\subsection{Arithmétique étendue et sa mesurabilité}
% ------------------------------------------------------------------

\begin{definition}[Opérations sur $\Rpb$]\label{def:arith}
On prolonge l'addition et la multiplication à $\Rpb$ par :
\[
a+(+\infty)=(+\infty)+a=+\infty\quad(a\in\Rpb),
\]
\[
a\cdot(+\infty)=(+\infty)\cdot a=+\infty\ (a\in\,]0,+\infty]),\qquad
\boxed{\,0\cdot(+\infty)=(+\infty)\cdot0=0\,.}
\]
La convention $0\cdot\infty=0$ est celle de la théorie de la mesure (elle rend
l'intégrale $\int 0\,d\mu=0$ même sur un ensemble de mesure infinie).
\end{definition}

\begin{remark}[Absence de forme indéterminée additive]
Sur $\Rpb$ \emph{il n'y a pas de forme $\infty-\infty$}, puisqu'on ne
soustrait pas : l'addition est partout définie et l'on va voir qu'elle est même
\emph{continue}. C'est un avantage propre à $\Rpb$ par rapport à la droite
achevée $\Rb=[-\infty,+\infty]$, où $(+\infty)+(-\infty)$ n'est pas défini.
\end{remark}

\begin{proposition}[Continuité de l'addition]\label{prop:add-continue}
L'addition $s:\Rpb\times\Rpb\to\Rpb$, $s(x,y)=x+y$, est continue (pour la
topologie produit). En particulier elle est borélienne.
\end{proposition}

\begin{proof}
On vérifie la continuité séquentielle, licite car $\Rpb\times\Rpb$ est
métrisable (produit de deux métriques compacts). Soit $(x_n,y_n)\to(x,y)$.
\begin{itemize}[nosep]
\item Si $x,y<+\infty$ : au voisinage de $(x,y)$, $x_n,y_n$ sont finis et
$x_n+y_n\to x+y$ par continuité de l'addition réelle (et $\varphi$ étant un
homéomorphisme, la convergence dans $\Rpb$ des suites finies équivaut à la
convergence usuelle).
\item Si $x=+\infty$ (le cas $y=+\infty$ est symétrique) : alors $x_n\to+\infty$,
i.e.\ pour tout $M$, $x_n>M$ à partir d'un certain rang. Comme $y_n\ge0$, on a
$x_n+y_n\ge x_n>M$ à partir de ce rang, donc $x_n+y_n\to+\infty=x+y$.
\end{itemize}
Dans tous les cas $s(x_n,y_n)\to s(x,y)$. Continue $\Rightarrow$ borélienne
(l'image réciproque d'un ouvert est un ouvert, donc un borélien).
\end{proof}

\begin{proposition}[Mesurabilité de la multiplication]\label{prop:mult-mesurable}
La multiplication $m:\Rpb\times\Rpb\to\Rpb$, $m(x,y)=x\cdot y$ (avec
$0\cdot\infty=0$), est \emph{borélienne}, mais \emph{discontinue} aux points
$(0,+\infty)$ et $(+\infty,0)$.
\end{proposition}

\begin{proof}
\emph{Discontinuité en $(0,+\infty)$.} Prenons $x_n=1/n\to0$ et
$y_n=+\infty$ constant : $x_n y_n=+\infty\not\to0=m(0,+\infty)$. Donc $m$ n'est
pas continue en $(0,+\infty)$ ; idem en $(+\infty,0)$ par symétrie.

\emph{Borélianité.} Sur l'ouvert
$U=\Rpb\times\Rpb\setminus\{(0,+\infty),(+\infty,0)\}$, la fonction $m$ est
continue. En effet, sur $U$ on n'a jamais simultanément un facteur nul et un
facteur infini ; la vérification par suites est identique à celle de
l'addition (si un facteur tend vers $+\infty$ et l'autre reste minoré par un
$\eps>0$, le produit tend vers $+\infty$ ; si les deux sont finis, produit
réel usuel). Ainsi $m|_U$ est continue donc borélienne sur l'ouvert $U$.

Il reste les deux points exceptionnels. L'ensemble
$E=\{(0,+\infty),(+\infty,0)\}$ est fermé (donc borélien) et $m$ y vaut $0$.
Pour tout borélien $B\subseteq\Rpb$,
\[
m^{-1}(B)=\big(m^{-1}(B)\cap U\big)\ \cup\ \big(m^{-1}(B)\cap E\big).
\]
Le premier terme est borélien car $m|_U$ est borélienne (continue sur l'ouvert
$U$) et $U$ est borélien. Le second est $E\cap m^{-1}(B)$, qui vaut $E$ si
$0\in B$ et $\varnothing$ sinon ; dans les deux cas c'est un borélien. Donc
$m^{-1}(B)\in\B(\Rpb\times\Rpb)$ et $m$ est borélienne.
\end{proof}

\begin{theorem}[Somme et produit de fonctions mesurables]\label{thm:somme-produit}
Soient $f,g:\Omega\to\Rpb$ mesurables. Alors $f+g$ et $f\cdot g$ (avec les
conventions de la Définition~\ref{def:arith}) sont mesurables.
\end{theorem}

\begin{proof}
L'application $\Phi:\Omega\to\Rpb\times\Rpb$, $\Phi(\omega)=(f(\omega),g(\omega))$,
est mesurable pour la tribu produit $\B(\Rpb)\otimes\B(\Rpb)$, car ses
composantes le sont. Or, $\Rpb$ étant à base dénombrable, on a l'égalité
\[
\B(\Rpb\times\Rpb)=\B(\Rpb)\otimes\B(\Rpb)
\]
(la tribu borélienne du produit d'espaces à base dénombrable est le produit
des tribus boréliennes : voir Lemme~\ref{lem:borel-produit}). Donc $\Phi$ est
$(\F,\B(\Rpb\times\Rpb))$-mesurable.

Comme $s(x,y)=x+y$ est borélienne (Proposition~\ref{prop:add-continue}) et
$m(x,y)=xy$ est borélienne (Proposition~\ref{prop:mult-mesurable}), les
composées
\[
f+g=s\circ\Phi,\qquad f\cdot g=m\circ\Phi
\]
sont mesurables (composée d'une application mesurable et d'une application
borélienne).
\end{proof}

\begin{lemma}[Borélien d'un produit à base dénombrable]\label{lem:borel-produit}
Si $X,Y$ sont des espaces topologiques à base dénombrable, alors
$\B(X\times Y)=\B(X)\otimes\B(Y)$.
\end{lemma}

\begin{proof}
\emph{$\supseteq$ :} les projections $\pi_X,\pi_Y$ sont continues donc
boréliennes, ce qui donne $\B(X)\otimes\B(Y)\subseteq\B(X\times Y)$ (la tribu
produit est la plus petite rendant les projections mesurables).

\emph{$\subseteq$ :} soient $\mathcal{U},\mathcal{V}$ des bases dénombrables de
$X,Y$. Les pavés $U\times V$ ($U\in\mathcal{U}$, $V\in\mathcal{V}$) forment une
base dénombrable de $X\times Y$. Tout ouvert $W\subseteq X\times Y$ est donc
réunion \emph{dénombrable} de tels pavés :
$W=\bigcup_k U_k\times V_k$. Or $U_k\times V_k\in\B(X)\otimes\B(Y)$, et une
réunion dénombrable reste dans la tribu produit, donc
$W\in\B(X)\otimes\B(Y)$. Comme les ouverts engendrent $\B(X\times Y)$, on
conclut $\B(X\times Y)\subseteq\B(X)\otimes\B(Y)$.
\end{proof}

\begin{remark}
L'hypothèse « base dénombrable » est essentielle : pour des espaces trop gros,
$\B(X\times Y)$ peut être strictement plus grande que $\B(X)\otimes\B(Y)$. Ici
$\Rpb$ est à base dénombrable (Proposition~\ref{prop:base-denombrable}), tout
va bien.
\end{remark}

% ------------------------------------------------------------------
\subsection{Principe de transfert par homéomorphisme}
% ------------------------------------------------------------------

\begin{theorem}[Transfert via $\varphi$]\label{thm:transfert}
Soit $\varphi:\Rpb\to[0,1]$ l'homéomorphisme de la Définition~\ref{def:phi}.
Pour toute application $g:\Omega\to\Rpb$,
\[
g\text{ est }(\F,\B(\Rpb))\text{-mesurable}
\iff
\varphi\circ g\text{ est }(\F,\B([0,1]))\text{-mesurable}.
\]
Ainsi toute question de mesurabilité à valeurs dans $\Rpb$ se ramène à une
question de mesurabilité à valeurs réelles bornées dans $[0,1]$.
\end{theorem}

\begin{proof}
$\varphi$ est un homéomorphisme (Proposition~\ref{prop:coinc-metrique}), donc
un isomorphisme borélien : $\varphi$ et $\varphi^{-1}$ sont boréliennes.

$(\Rightarrow)$ Si $g$ est mesurable, alors $\varphi\circ g$ est mesurable comme
composée d'une application mesurable $g$ et d'une application borélienne
$\varphi$.

$(\Leftarrow)$ Si $\varphi\circ g$ est mesurable, alors
$g=\varphi^{-1}\circ(\varphi\circ g)$ est mesurable comme composée de
$\varphi\circ g$ (mesurable) et de $\varphi^{-1}$ (borélienne).
\end{proof}

\begin{corollary}
Les Théorèmes~\ref{thm:supinf},~\ref{thm:limsup} et~\ref{thm:somme-produit}
peuvent alternativement se déduire des énoncés analogues sur $[0,1]$ (voire sur
$\R$) par transfert, à condition de traduire correctement l'arithmétique
étendue via $\varphi$.
\end{corollary}

% ==================================================================
\section{Synthèse}
% ==================================================================

Le tableau suivant récapitule la double structure de $\Rpb$ et ses
conséquences.

\begin{center}
\renewcommand{\arraystretch}{1.35}
\begin{tabular}{@{}p{5.2cm}p{9.2cm}@{}}
\hline
\textbf{Structure} & \textbf{Conséquence}\\
\hline
Treillis complet (Prop.~\ref{prop:treillis}) &
$\sup,\inf$ d'une famille quelconque existent dans $\Rpb$ ;
$\limsup,\liminf$ toujours définis.\\
Topologie de l'ordre $=$ Alexandroff $=$ métrique via $\varphi$
(Prop.~\ref{prop:coinc-metrique},~\ref{prop:alex-ordre}) &
compactification canonique ; $\Rpb\cong[0,1]$ compact métrisable.\\
Base dénombrable (Prop.~\ref{prop:base-denombrable}) &
espace borélien standard ; $\B(\Rpb\times\Rpb)=\B(\Rpb)^{\otimes2}$.\\
$\{+\infty\}$ fermé $G_\delta$ (Lem.~\ref{lem:infty-borel}) &
caractérisation par trace : $B$ ou $B\cup\{+\infty\}$, $B\in\B(\Rp)$
(Th.~\ref{thm:trace}).\\
Générateurs $(a,+\infty]$ (Th.~\ref{thm:generateurs}) &
critère $\{f>a\}\in\F$ pour la mesurabilité (Cor.~\ref{cor:criteres}).\\
Addition continue, produit borélien
(Prop.~\ref{prop:add-continue},~\ref{prop:mult-mesurable}) &
$f+g$, $fg$ mesurables (Th.~\ref{thm:somme-produit}).\\
Homéomorphisme $\varphi$ (Th.~\ref{thm:transfert}) &
transfert de toute la mesurabilité vers $[0,1]$.\\
\hline
\end{tabular}
\end{center}

\medskip
\noindent
\textbf{En une phrase.} Le point $+\infty$ transforme $\Rp$ en un treillis
complet \emph{et} en un compact métrisable (le segment) ; la topologie
sous-jacente est canonique (ordre $=$ Alexandroff $=$ métrique) ; la tribu
borélienne se lit comme la trace « $B$ ou $B\cup\{+\infty\}$ » ; et c'est cette
double structure — ordre-complet d'un côté, compacte de l'autre — qui rend
$\sup$, $\inf$, $\limsup$, $\liminf$, ainsi que $+$ et $\times$, mesurables et
partout définis.

\end{document}
